% GENERATED FILE. Edit canonical lessons, then run npm run generate:books.
% canonical-source-hash: fnv1a64:35c88dff9cd37a3e
% canonical-lessons: SA-C256-prasiksanam, SA-C256-sammelanam, SA-C256-nirnayah, SA-C256-vicarah, SA-C256-matam

\chapter{Training, Conference, Decision, Thought, Opinion}
\label{ch:sa-training-conference-decision-thought-opinion}
\hlchaptermodality{256}

\begin{chapteropening}
\textbf{By the end of this chapter:} \emph{I can say training, a conference, a decision, a thought and an opinion.}

Complete the last lesson of chapter 256: I can say training, a conference, a decision, a thought and an opinion.
\end{chapteropening}

\section[\texorpdfstring{\sk{प्रशिक्षणम्} (praśikṣaṇam)}{प्रशिक्षणम् (praśikṣaṇam)}]{\sk{प्रशिक्षणम्} (praśikṣaṇam) — training}

\label{lesson:SA-C256-prasiksanam}



\begin{quote}
Before the new one: say the Sanskrit for homework, then the Sanskrit for an answer.
\end{quote}

\subsection*{You'll want to know: \sk{प्रशिक्षणम्}}
\textbf{\sk{प्रशिक्षणम्}} — \emph{praśikṣaṇam} — \textquotedblleft{}training\textquotedblright{}.

\begin{cousinweb}
\emph{praśikṣaṇa} is built on \emph{pra-} and \emph{śikṣaṇa}, \textquotedblleft{}learning; teaching, instruction\textquotedblright{} (MW), from \emph{śikṣ}, the desiderative of √śak, \textquotedblleft{}to be able\textquotedblright{}. MW does not list \emph{praśikṣaṇa}, so the word for training is a modern coinage.
\end{cousinweb}

\begin{grammarlens}[title={What you've built}]
One more word for this chapter.
\end{grammarlens}

\subsection*{Your turn}
\begin{itemize}
\raggedright
  \item \emph{Say it:} \emph{praśikṣaṇam}
  \item \emph{Say it:} \emph{praśikṣaṇam}, once more
  \item \emph{Say it:} say \emph{praśikṣaṇam}
  \item \emph{Recall:} say the Sanskrit for homework, then the Sanskrit for an answer, then say \emph{praśikṣaṇam} again
\end{itemize}

\begin{tcolorbox}[breakable,colback=teal!4,colframe=teal!35!black,title={Before you move on}]
What does \sk{प्रशिक्षणम्} mean? (\textquotedblleft{}Training\textquotedblright{}.) Say it once more.
\end{tcolorbox}
\section[\texorpdfstring{\sk{सम्मेलनम्} (sammelanam)}{सम्मेलनम् (sammelanam)}]{\sk{सम्मेलनम्} (sammelanam) — a conference}

\label{lesson:SA-C256-sammelanam}



\begin{quote}
Before the new one: say the Sanskrit for an answer, then the Sanskrit for training.
\end{quote}

\subsection*{You'll want to know: \sk{सम्मेलनम्}}
\textbf{\sk{सम्मेलनम्}} — \emph{sammelanam} — \textquotedblleft{}a conference\textquotedblright{}.

\begin{cousinweb}
\emph{sammelana} is built on \emph{sam}, \textquotedblleft{}together\textquotedblright{}, and the root √mil, \textquotedblleft{}to meet\textquotedblright{} (MW): \textquotedblleft{}mingling or meeting together\textquotedblright{}. The same root gives \emph{melaka}, \textquotedblleft{}fair\textquotedblright{}, which you have already met, and Hindi \emph{milnā}, \textquotedblleft{}to meet\textquotedblright{} (CDIAL 10133).
\end{cousinweb}

\begin{grammarlens}[title={What you've built}]
One more word for this chapter.
\end{grammarlens}

\subsection*{Your turn}
\begin{itemize}
\raggedright
  \item \emph{Say it:} \emph{sammelanam}
  \item \emph{Say it:} \emph{sammelanam}, once more
  \item \emph{Say it:} say \emph{sammelanam}
  \item \emph{Recall:} say the Sanskrit for an answer, then the Sanskrit for training, then say \emph{sammelanam} again
\end{itemize}

\begin{tcolorbox}[breakable,colback=teal!4,colframe=teal!35!black,title={Before you move on}]
What does \sk{सम्मेलनम्} mean? (\textquotedblleft{}A conference\textquotedblright{}.) Say it once more.
\end{tcolorbox}
\section[\texorpdfstring{\sk{निर्णयः} (nirṇayaḥ)}{निर्णयः (nirṇayaḥ)}]{\sk{निर्णयः} (nirṇayaḥ) — a decision}

\label{lesson:SA-C256-nirnayah}



\begin{quote}
Before the new one: say the Sanskrit for training, then the Sanskrit for a conference.
\end{quote}

\subsection*{You'll want to know: \sk{निर्णयः}}
\textbf{\sk{निर्णयः}} — \emph{nirṇayaḥ} — \textquotedblleft{}a decision\textquotedblright{}.

\begin{cousinweb}
\emph{nirṇaya} is built on \emph{nis}, \textquotedblleft{}out\textquotedblright{}, and the root √nī, \textquotedblleft{}to lead\textquotedblright{} (MW). The verb means \textquotedblleft{}to carry off\textquotedblright{} and then \textquotedblleft{}to find out, settle, decide\textquotedblright{}, and the noun means \textquotedblleft{}complete ascertainment, decision\textquotedblright{}.
\end{cousinweb}

\begin{grammarlens}[title={What you've built}]
One more word for this chapter.
\end{grammarlens}

\subsection*{Your turn}
\begin{itemize}
\raggedright
  \item \emph{Say it:} \emph{nirṇayaḥ}
  \item \emph{Say it:} \emph{nirṇayaḥ}, once more
  \item \emph{Say it:} say \emph{nirṇayaḥ}
  \item \emph{Recall:} say the Sanskrit for training, then the Sanskrit for a conference, then say \emph{nirṇayaḥ} again
\end{itemize}

\begin{tcolorbox}[breakable,colback=teal!4,colframe=teal!35!black,title={Before you move on}]
What does \sk{निर्णयः} mean? (\textquotedblleft{}A decision\textquotedblright{}.) Say it once more.
\end{tcolorbox}
\section[\texorpdfstring{\sk{विचारः} (vicāraḥ)}{विचारः (vicāraḥ)}]{\sk{विचारः} (vicāraḥ) — a thought, an idea}

\label{lesson:SA-C256-vicarah}



\begin{quote}
Before the new one: say the Sanskrit for a conference, then the Sanskrit for a decision.
\end{quote}

\subsection*{You'll want to know: \sk{विचारः}}
\textbf{\sk{विचारः}} — \emph{vicāraḥ} — \textquotedblleft{}a thought, an idea\textquotedblright{}.

\begin{cousinweb}
\emph{vicāra} is built on \emph{vi-}, \textquotedblleft{}apart, about\textquotedblright{}, and the root √car, \textquotedblleft{}to move\textquotedblright{} (MW): first \textquotedblleft{}mode of acting\textquotedblright{}, then \textquotedblleft{}pondering, deliberation\textquotedblright{} (CDIAL 11683). Thought is the mind moving to and fro.
\end{cousinweb}

\begin{grammarlens}[title={What you've built}]
One more word for this chapter.
\end{grammarlens}

\subsection*{Your turn}
\begin{itemize}
\raggedright
  \item \emph{Say it:} \emph{vicāraḥ}
  \item \emph{Say it:} \emph{vicāraḥ}, once more
  \item \emph{Say it:} say \emph{vicāraḥ}
  \item \emph{Recall:} say the Sanskrit for a conference, then the Sanskrit for a decision, then say \emph{vicāraḥ} again
\end{itemize}

\begin{tcolorbox}[breakable,colback=teal!4,colframe=teal!35!black,title={Before you move on}]
What does \sk{विचारः} mean? (\textquotedblleft{}A thought, an idea\textquotedblright{}.) Say it once more.
\end{tcolorbox}
\section[\texorpdfstring{\sk{मतम्} (matam)}{मतम् (matam)}]{\sk{मतम्} (matam) — an opinion}

\label{lesson:SA-C256-matam}



\begin{quote}
Before the new one: say the Sanskrit for a decision, then the Sanskrit for a thought.
\end{quote}

\subsection*{You'll want to know: \sk{मतम्}}
\textbf{\sk{मतम्}} — \emph{matam} — \textquotedblleft{}an opinion\textquotedblright{}.

\begin{cousinweb}
\emph{mata} is the past participle of the root √man, \textquotedblleft{}to think\textquotedblright{} (MW): \textquotedblleft{}what is thought\textquotedblright{}, an opinion. MW compares √man with Latin \emph{meminisse}, \textquotedblleft{}to remember\textquotedblright{}, and \emph{monēre}, \textquotedblleft{}to remind\textquotedblright{}. The same root gives \emph{manas}, \textquotedblleft{}mind\textquotedblright{}.
\end{cousinweb}

\begin{grammarlens}[title={What you've built}]
One more word for this chapter.
\end{grammarlens}

\subsection*{Your turn}
\begin{itemize}
\raggedright
  \item \emph{Say it:} \emph{matam}
  \item \emph{Say it:} \emph{matam}, once more
  \item \emph{Say it:} say \emph{matam}
  \item \emph{Recall:} say the Sanskrit for a decision, then the Sanskrit for a thought, then say \emph{matam} again
\end{itemize}

\begin{tcolorbox}[breakable,colback=teal!4,colframe=teal!35!black,title={Before you move on}]
What does \sk{मतम्} mean? (\textquotedblleft{}An opinion\textquotedblright{}.) Say it once more.
\end{tcolorbox}
